% PREAMPLE

\documentclass[a4paper, 14pt]{article}

\usepackage[
    inner=10mm,
    outer=10mm,
    top=5mm,
    bottom=10mm
]{geometry}
\usepackage{xcolor}
\usepackage{tikz}
\usepackage{hyperref}
\usepackage{ulem}
\usepackage{float}
\usepackage{fontawesome}

\renewcommand{\familydefault}{\sfdefault}

\title{Thomas Møller Jensen \\
\tiny MSc Software Engineering / BSc Computer Engineering \\
\includegraphics[width=.08\textwidth]{cache/261020_aau-logo/__AAU_LEFT_RGB_UK.png}
\small / \href{https://github.com/Mast3rwaf1z}{\includegraphics[width=.02\textwidth]{cache/github-mark/github-mark.png} Mast3rwaf1z}
/ \href{https://about.skademaskinen.win}{Website}}
\date{}
\author{}
\pagenumbering{gobble}

\definecolor{myorange}{HTML}{dd2200}
\newcommand{\myitem}[1]{\item[\color{myorange}$\circ$]#1}

\newcommand{\cvsection}[2]{
    \begin{minipage}[t]{.25\textwidth}
        [ \textbf{#1} ]:
    \end{minipage}
    \vrule
    \hspace{2mm}
    \begin{minipage}[t]{.70\textwidth}
        #2
    \end{minipage}
    \vspace{.25cm}
}

\newcommand{\newcvsection}[2]{
    \noindent\huge \textbf{#1} \color{myorange}\hrulefill\color{black}
    
    \normalsize
    #2    
    \vspace{.5cm}
}


\begin{document}
    % FORMALIA
    \begin{minipage}{.75\textwidth}
    \maketitle
    \end{minipage}
    \begin{minipage}{.25\textwidth}
        \begin{tikzpicture}
            \clip (0,0) circle (1.5cm);
            \node[anchor=center] at (0,0) {\includegraphics[width=.8\textwidth]{files/pfp.jpg}};
        \end{tikzpicture}
    \end{minipage}

    % MAIN

    %\cvsection{Me}{
    %    Hi, as of writing my CV, I am an 8th semester student in Software Engineering at Aalborg University. 
    %    I have a bachelor in Computer Engineering. 
    %    My personal interests include; gaming, D\&D, programming and Linux.
    %    In terms of programming I take particular interest in functional programming, but anything that can enhance my experience as a Linux user is also interesting to me.
    %    As a person who plays video games in Linux, it also provides challenges, so i have experience with tweaking various graphical drivers and layers in the graphical software stack in Linux.
    %}

    \newcvsection{Education}{
        \noindent\textbf{BSc Computer Engineering} \hfill \color{gray} Sep. 2020 - Jun. 2023
        
        \color{black} Aalborg University - Department of Electronic Systems

        \noindent\textbf{MSc Software Engineering} \hfill \color{gray} Sep. 2023 - Jun. 2025

        \color{black} Aalborg University - Department of Computer Science
    }

    \newcvsection{Skills}{
        \begin{table}[H]
            \begin{tabular}{rl}
                \textbf{Programming} & Python, C, C++, Nix, Java, Haskell\\
                \textbf{OS} & Arch, NixOS, anything Linux... \\
                \textbf{Build tools} & Cmake, Gradle, Maven \\
                \textbf{Languages} & English, Danish
            \end{tabular}
        \end{table}
    }

    \newcvsection{Significant Projects}{
        \begin{itemize}
            \myitem{
                \textbf{Bachelor Project: } An interesting research project where we made a series of tests for a testbed and actually acquired some interesting results. 
                In this project I learned a lot about computer networking and tools you can use to interface with the networking functionality in the Linux kernel.
            }
            \myitem{
                \textbf{Skademaskinen: } My server, in conjunction with a friend from Aalborg University we have used various computer parts to create a server which hosts a variety of software we have written.
                Available on \href{https://github.com/Skademaskinen}{Github}.
                \begin{itemize}
                    \myitem{
                        Two of the three Discord bots are written by me in Java using the dv8tion JDA library. One of these is thoroughly tested using GitHub Actions.
                    }
                    \myitem{
                        My website has a backend server hosting a REST API, handling authorization and my home network, as well as a method for me to update the site without changing the source code.
                    }
                \end{itemize}
            }

            \myitem{
                \textbf{ASCK: } As a member of Aalborg Supercomputer Klub I have contributed scripts for the clusters (runners and a htop dashboard using tmux), participation in meetings and competitions. 
                I have experience working with MPI (mpich) from these competitions, but also knowledge about our setups for HPL and HPCG.
            }
        \end{itemize}
    }
    
    \newcvsection{Work Experience}{
        \begin{itemize}
            \myitem{
                \textbf{DEIS: } During my bachelor studies, I worked as a student programmer at DEIS, a research unit of Department of Computer Science at Aalborg University. 
                This job was focused on making a GUI solution for a framework for a course in Distributed Systems.
            }
                
            \myitem{
                \textbf{F-Klubben: } I have worked as a volunteer at F-klubben, an association for students and employees at Department of Computer Science. 
                As a member of the association my role has primarily been as a contributer to their software, most notably working on \href{https://github.com/f-klubben/fnix}{fnix}.
            }
        \end{itemize}
    }

    \newcvsection{Competitions}{
        \begin{itemize}
            \myitem{
                \textbf{SBCC24: } I was one of four competitors on the team from ASCK in the Single Board Cluster Competition 2024.
                Along with AAU, there were also teams from Kansas, San Diego and Texas universities.
                In this competition we ran HPL, HPCG, a rubiks cube solving program (twophase), a remote computation service called BOINC and some standard benchmarks called SU2.
                Out of all these programs, ASCK got first place in all but one, winning us the overall competition.
            }
        \end{itemize}
    }

    % FOOTER
    \vspace*{\fill}
    \tiny Written in \LaTeX, Source is available at: \href{https://github.com/Mast3rwaf1z/CV}{https://github.com/Mast3rwaf1z/CV}
\end{document}